% Lösungen Einheit 1 von 3 – Der Quotient (zusammenbau v0.1)
\einheitenkopf{Einheit 1 von 3 · Der Quotient}

\erg{8}{a) die Fußball spielenden Mitglieder \quad b) $\frac{0{,}16}{0{,}4} = 0{,}4$ \quad c) $\frac{36}{90} = 0{,}4$ \quad d) Mädchen: $\frac{15}{50} = 0{,}3$; Jungen: $\frac{14}{35} = 0{,}4$; unter den Jungen ist der Anteil größer, obwohl mehr Mädchen eine Brille tragen \quad e) $P_B(A) = \frac{0{,}12}{0{,}3} = 0{,}4$ und $P_A(B) = \frac{0{,}12}{0{,}4} = 0{,}3$; das Feld ist dasselbe, geteilt wird aber einmal durch den Anteil der Stadtbewohner, einmal durch den Anteil der Allergiker}
\erg{9}{a) Fehler: Ben teilt durch den Anteil der Frühstücksgäste, die Bedingung ist aber „Geschäftsreisender“; richtig $\frac{0{,}14}{0{,}28} = 0{,}5$ \quad b) Die Bedingung schränkt auf eine Gruppe ein; gefragt ist ein Anteil innerhalb dieser Gruppe, also der Teil der Gruppe durch die ganze Gruppe \quad c) Bus: $\frac{0{,}22}{0{,}55} = 0{,}4$; übrige: $\frac{0{,}08}{0{,}45} \approx 0{,}178$; ja, Busfahrer kommen deutlich häufiger zu spät}
